\documentclass{beamer}

\usetheme{Madrid}

% Custom footer
\setbeamertemplate{footline}{
  \leavevmode%
  \hbox{%
  \begin{beamercolorbox}[wd=.333333\paperwidth,ht=2.25ex,dp=1ex,center]{author in head/foot}%
    de Mel, Raj, Wasif, Brake
\end{beamercolorbox}%
  \begin{beamercolorbox}[wd=.333333\paperwidth,ht=2.25ex,dp=1ex,center]{title in head/foot}%
    COMP3242/COMP6242 - Deep Learning
  \end{beamercolorbox}%
  \begin{beamercolorbox}[wd=.333333\paperwidth,ht=2.25ex,dp=1ex,right]{date in head/foot}%
    31 May 2026 \hspace*{2em} \insertframenumber{} / \inserttotalframenumber \hspace*{2ex}
  \end{beamercolorbox}}%
  \vskip0pt%
}

% Title page info
\title{Multi-Backbone Ensemble with LoRA Adaptation for
Multi-Label Plant Species Identification in Vegetation
Plots}
\author{Manindra de Mel, Arjun Raj, Razeen Wasif, Will Brake}
\institute{\textbf{Australian National University}}
\date{31 May 2026}

\begin{document}

\frame{\titlepage}

\begin{frame}{Introduction - Task and Motivation}

\begin{itemize}
    \item Vegetation surveys are \textbf{costly} and \textbf{hard to scale} - manual quadrat-based identification requires expert botanists on-site
\end{itemize}
\begin{itemize}
    \item PlantCLEF 2026 \textbf{automates} this - multi-label classification of all plant species from a single quadrat photo across ~7,800 species
\end{itemize}
\begin{itemize}
    \item Evaluated using \textbf{macro-averaged F1} - computed per sample and aggregated across transects
\end{itemize}
$$\text{Score} = \frac{1}{N} \sum_{i=1}^{N} \left( \frac{1}{T_i} \sum_{j=1}^{T_i} \text{F1}^j \right)$$Where $N$ is the total number of transects. $T_i$ is the number of quadrats in transect $i$. $\text{F1}^j$ is the macro-averaged F1 score for test image $j$. 
    
\end{frame}

\begin{frame}{Introduction - Challenges}

\begin{itemize}
    \item \textbf{Domain shift} between training and test data - models trained on isolated single-species images struggle with dense, complex quadrat scenes where multiple species overlap and obstruct one another
\end{itemize}
\begin{figure}
    \centering
    \includegraphics[width=0.45\linewidth]{DomainShift.png}
    \caption{Illustration of the visual discrepancy between test and training sets. (Source: PlantCLEF 2026 Challenge, Kaggle)}
\end{figure}
\begin{itemize}
    \item Long-tail species distribution - common species have thousands of training images while rare species may have fewer than ten, causing models to \textbf{over-predict} common classes and yielding \textbf{low} macro-F1
\end{itemize}
    
\end{frame}

\begin{frame}{Method --- Pipeline Overview}
    \begin{center}
        \includegraphics[width=\textwidth]{pipeline_diagram_.pdf}
    \end{center}
\end{frame}

\begin{frame}{Experiments}
    \begin{center}
        \includegraphics[width=\textwidth]{}
\begin{figure}
        \centering
        \includegraphics[width=1\linewidth]{score_progression_madrid.pdf}
    \end{figure}
        \end{center}
    
\end{frame}

\begin{frame}{Key Findings}
  \begin{itemize}
    \item Fine-tuning only the top four transformer blocks
      outperformed a full fine-tune by 20 F1 points, preserving
      the model's pretrained biological knowledge.
    \item Validation accuracy increased while leaderboard
      performance stayed the same - the single-plant to multi-species
      domain gap, is the main constraint.
    \item Additional priors produced no measurable gain as genus
      co-occurrence and seasonal timing were both already
      found by the model during training.
  \end{itemize}
\end{frame}

\begin{frame}{Conclusion and Further Work}
  \begin{itemize}
    \item Best result: \textbf{0.412 public / 0.403 private
      Macro-F1} from a single BioCLIP~2.5 checkpoint, without
      ensembling.
    \item More potential gains from improving data  then better architecture
      - training on synthetic multi-species scenes that match
      the test distribution.
    \item New signal could come from \textbf{outside
      the image}: location, observation date, and per-class
      threshold calibration.
  \end{itemize}
\end{frame}


\end{document}

